% !TeX program = xelatex
% 由生成课件.py生成；改题目请改共享题库，改选题请改课堂选题.json。
\documentclass[aspectratio=169,10pt]{beamer}
\usepackage[UTF8,fontset=fandol]{ctex}
\IfFontExistsTF{Microsoft YaHei}{\setCJKsansfont{Microsoft YaHei}}{}
\usetheme{GaoZhongYuan}
\usepackage{booktabs}
\setlength{\parskip}{.6em}
\newcommand{\ans}[1]{{\color{red!70!black}#1}}
\author{深圳中学高中园}\date{}

\title{对数的概念：全量备课题库}
\begin{document}\campusmaketitle
% 题库编号：LOG-001
\begin{campuspractice}{题库 LOG-001}
将下式在指数式与对数式之间互化：\par
$5^4=625$
\end{campuspractice}

\ifdefined\TeacherEdition
\begin{frame}[t]{LOG-001：参考解答}
$\log_5625=4$
\end{frame}
\fi

% 题库编号：LOG-002
\begin{campuspractice}{题库 LOG-002}
将下式在指数式与对数式之间互化：\par
$2^{-6}=\frac{1}{64}$
\end{campuspractice}

\ifdefined\TeacherEdition
\begin{frame}[t]{LOG-002：参考解答}
$\log_2\frac1{64}=-6$
\end{frame}
\fi

% 题库编号：LOG-003
\begin{campuspractice}{题库 LOG-003}
将下式在指数式与对数式之间互化：\par
$(\frac{1}{3})^m=5.73$
\end{campuspractice}

\ifdefined\TeacherEdition
\begin{frame}[t]{LOG-003：参考解答}
$\log_{1/3}5.73=m$
\end{frame}
\fi

% 题库编号：LOG-004
\begin{campuspractice}{题库 LOG-004}
将下式在指数式与对数式之间互化：\par
$\log_{\frac{1}{2}}16=-4$
\end{campuspractice}

\ifdefined\TeacherEdition
\begin{frame}[t]{LOG-004：参考解答}
$(\frac12)^{-4}=16$
\end{frame}
\fi

% 题库编号：LOG-005
\begin{campuspractice}{题库 LOG-005}
将下式在指数式与对数式之间互化：\par
$\lg0.01=-2$
\end{campuspractice}

\ifdefined\TeacherEdition
\begin{frame}[t]{LOG-005：参考解答}
$10^{-2}=0.01$
\end{frame}
\fi

% 题库编号：LOG-006
\begin{campuspractice}{题库 LOG-006}
将下式在指数式与对数式之间互化：\par
$\ln10\approx2.303$
\end{campuspractice}

\ifdefined\TeacherEdition
\begin{frame}[t]{LOG-006：参考解答}
$e^{2.303}\approx10$
\end{frame}
\fi

% 题库编号：LOG-007
\begin{campuspractice}{题库 LOG-007}
求下式中的 $x$：\par
$\log_{64}x=\frac{-2}{3}$
\end{campuspractice}

\ifdefined\TeacherEdition
\begin{frame}[t]{LOG-007：参考解答}
$x=\frac1{16}$
\end{frame}
\fi

% 题库编号：LOG-008
\begin{campuspractice}{题库 LOG-008}
求下式中的 $x$：\par
$\log_x8=6$
\end{campuspractice}

\ifdefined\TeacherEdition
\begin{frame}[t]{LOG-008：参考解答}
$x=\sqrt2$
\end{frame}
\fi

% 题库编号：LOG-009
\begin{campuspractice}{题库 LOG-009}
求下式中的 $x$：\par
$\lg100=x$
\end{campuspractice}

\ifdefined\TeacherEdition
\begin{frame}[t]{LOG-009：参考解答}
$x=2$
\end{frame}
\fi

% 题库编号：LOG-010
\begin{campuspractice}{题库 LOG-010}
求下式中的 $x$：\par
$-\ln(e^2)=x$
\end{campuspractice}

\ifdefined\TeacherEdition
\begin{frame}[t]{LOG-010：参考解答}
$x=-2$
\end{frame}
\fi

% 题库编号：LOG-011
\begin{campuspractice}{题库 LOG-011}
$\log_3\frac{2x-1}{5}=0$
\end{campuspractice}

\ifdefined\TeacherEdition
\begin{frame}[t]{LOG-011：参考解答}
$x=3$
\end{frame}
\fi

% 题库编号：LOG-012
\begin{campuspractice}{题库 LOG-012}
$6^{\log_6(5x+1)}=36$
\end{campuspractice}

\ifdefined\TeacherEdition
\begin{frame}[t]{LOG-012：参考解答}
$x=7$
\end{frame}
\fi

% 题库编号：LOG-013
\begin{campuspractice}{题库 LOG-013}
计算 $2^{\log_2(3+2)}$
\end{campuspractice}

\ifdefined\TeacherEdition
\begin{frame}[t]{LOG-013：参考解答}
$5$
\end{frame}
\fi

% 题库编号：LOG-014
\begin{campuspractice}{题库 LOG-014}
求下式有意义时 $x$ 的范围：\par
$\log_2(x-10)$
\end{campuspractice}

\ifdefined\TeacherEdition
\begin{frame}[t]{LOG-014：参考解答}
$x>10$
\end{frame}
\fi

% 题库编号：LOG-015
\begin{campuspractice}{题库 LOG-015}
求下式有意义时 $x$ 的范围：\par
$\log_{x-1}(x+2)$
\end{campuspractice}

\ifdefined\TeacherEdition
\begin{frame}[t]{LOG-015：参考解答}
$x\in(1,2)\cup(2,+\infty)$
\end{frame}
\fi

% 题库编号：LOG-016
\begin{campuspractice}{题库 LOG-016}
求下式中的 $x$：\par
$\log_2x=\frac{-1}{2}$
\end{campuspractice}

\ifdefined\TeacherEdition
\begin{frame}[t]{LOG-016：参考解答}
$x=\frac{\sqrt2}2$
\end{frame}
\fi

% 题库编号：LOG-017
\begin{campuspractice}{题库 LOG-017}
求下式中的 $x$：\par
$\log_x25=2$
\end{campuspractice}

\ifdefined\TeacherEdition
\begin{frame}[t]{LOG-017：参考解答}
$x=5$
\end{frame}
\fi

% 题库编号：LOG-018
\begin{campuspractice}{题库 LOG-018}
求下式中的 $x$：\par
$\log_5(x^2)=2$
\end{campuspractice}

\ifdefined\TeacherEdition
\begin{frame}[t]{LOG-018：参考解答}
$x=\pm5$
\end{frame}
\fi

% 题库编号：LOG-019
\begin{campuspractice}{题库 LOG-019}
将下式在指数式与对数式之间互化：\par
$5^4=625$
\end{campuspractice}

\ifdefined\TeacherEdition
\begin{frame}[t]{LOG-019：参考解答}
$\log_5625=4$
\end{frame}
\fi

% 题库编号：LOG-020
\begin{campuspractice}{题库 LOG-020}
将下式在指数式与对数式之间互化：\par
$2^{-6}=\frac{1}{64}$
\end{campuspractice}

\ifdefined\TeacherEdition
\begin{frame}[t]{LOG-020：参考解答}
$\log_2\frac1{64}=-6$
\end{frame}
\fi

% 题库编号：LOG-021
\begin{campuspractice}{题库 LOG-021}
将下式在指数式与对数式之间互化：\par
$4^x=10$
\end{campuspractice}

\ifdefined\TeacherEdition
\begin{frame}[t]{LOG-021：参考解答}
$\log_410=x$
\end{frame}
\fi

% 题库编号：LOG-022
\begin{campuspractice}{题库 LOG-022}
将下式在指数式与对数式之间互化：\par
$3^{-m/4}=9$
\end{campuspractice}

\ifdefined\TeacherEdition
\begin{frame}[t]{LOG-022：参考解答}
$\log_39=-m/4$
\end{frame}
\fi

% 题库编号：LOG-023
\begin{campuspractice}{题库 LOG-023}
将下式在指数式与对数式之间互化：\par
$\log_216=4$
\end{campuspractice}

\ifdefined\TeacherEdition
\begin{frame}[t]{LOG-023：参考解答}
$2^4=16$
\end{frame}
\fi

% 题库编号：LOG-024
\begin{campuspractice}{题库 LOG-024}
将下式在指数式与对数式之间互化：\par
$\log_3(\frac{1}{9})=-2$
\end{campuspractice}

\ifdefined\TeacherEdition
\begin{frame}[t]{LOG-024：参考解答}
$3^{-2}=1/9$
\end{frame}
\fi

% 题库编号：LOG-025
\begin{campuspractice}{题库 LOG-025}
将下式在指数式与对数式之间互化：\par
$\log_3(5+x)=2$
\end{campuspractice}

\ifdefined\TeacherEdition
\begin{frame}[t]{LOG-025：参考解答}
$3^2=5+x$
\end{frame}
\fi

% 题库编号：LOG-026
\begin{campuspractice}{题库 LOG-026}
将下式在指数式与对数式之间互化：\par
$\log_{a-1}(b+2)=3$
\end{campuspractice}

\ifdefined\TeacherEdition
\begin{frame}[t]{LOG-026：参考解答}
$(a-1)^3=b+2$，$a>1,a\ne2,b>-2$
\end{frame}
\fi

% 题库编号：LOG-027
\begin{campuspractice}{题库 LOG-027}
计算：\par
$\log_2(\frac{1}{16})$
\end{campuspractice}

\ifdefined\TeacherEdition
\begin{frame}[t]{LOG-027：参考解答}
$-4$
\end{frame}
\fi

% 题库编号：LOG-028
\begin{campuspractice}{题库 LOG-028}
计算：\par
$\log_{\frac{1}{2}}32$
\end{campuspractice}

\ifdefined\TeacherEdition
\begin{frame}[t]{LOG-028：参考解答}
$-5$
\end{frame}
\fi

% 题库编号：LOG-029
\begin{campuspractice}{题库 LOG-029}
计算：\par
$\log_{0.4}1$
\end{campuspractice}

\ifdefined\TeacherEdition
\begin{frame}[t]{LOG-029：参考解答}
$0$
\end{frame}
\fi

% 题库编号：LOG-030
\begin{campuspractice}{题库 LOG-030}
计算：\par
$\log_{2.5}6.25$
\end{campuspractice}

\ifdefined\TeacherEdition
\begin{frame}[t]{LOG-030：参考解答}
$2$
\end{frame}
\fi

% 题库编号：LOG-031
\begin{campuspractice}{题库 LOG-031}
计算：\par
$4^{\log_43}$
\end{campuspractice}

\ifdefined\TeacherEdition
\begin{frame}[t]{LOG-031：参考解答}
$3$
\end{frame}
\fi

% 题库编号：LOG-032
\begin{campuspractice}{题库 LOG-032}
计算：\par
$3^{\log_37+1}$
\end{campuspractice}

\ifdefined\TeacherEdition
\begin{frame}[t]{LOG-032：参考解答}
$21$
\end{frame}
\fi

% 题库编号：LOG-033
\begin{campuspractice}{题库 LOG-033}
求下式中的 $x$：\par
$\log_5x=3$
\end{campuspractice}

\ifdefined\TeacherEdition
\begin{frame}[t]{LOG-033：参考解答}
$x=125$
\end{frame}
\fi

% 题库编号：LOG-034
\begin{campuspractice}{题库 LOG-034}
求下式中的 $x$：\par
$\log_2(2x+1)=3$
\end{campuspractice}

\ifdefined\TeacherEdition
\begin{frame}[t]{LOG-034：参考解答}
$x=7/2$
\end{frame}
\fi

% 题库编号：LOG-035
\begin{campuspractice}{题库 LOG-035}
求下式中的 $x$：\par
$\log_x(\frac{1}{8})=3$
\end{campuspractice}

\ifdefined\TeacherEdition
\begin{frame}[t]{LOG-035：参考解答}
$x=1/2$
\end{frame}
\fi

% 题库编号：LOG-036
\begin{campuspractice}{题库 LOG-036}
求下式中的 $x$：\par
$\log_2(8^x)=-3$
\end{campuspractice}

\ifdefined\TeacherEdition
\begin{frame}[t]{LOG-036：参考解答}
$x=-1$
\end{frame}
\fi

% 题库编号：LOG-037
\begin{campuspractice}{题库 LOG-037}
$b=\log_{a-2}(5-a)$ 中，求 $a$ 的范围。
\end{campuspractice}

\ifdefined\TeacherEdition
\begin{frame}[t]{LOG-037：参考解答}
$a\in(2,3)\cup(3,5)$
\end{frame}
\fi

% 题库编号：LOG-038
\begin{campuspractice}{题库 LOG-038}
$\log_a(1-2a)$ 有意义时，求 $a$ 的范围。
\end{campuspractice}

\ifdefined\TeacherEdition
\begin{frame}[t]{LOG-038：参考解答}
$0<a<1/2$
\end{frame}
\fi

% 题库编号：LOG-039
\begin{campuspractice}{题库 LOG-039}
将下式互化：\par
$5^3=125$
\end{campuspractice}

\ifdefined\TeacherEdition
\begin{frame}[t]{LOG-039：参考解答}
$\log_5125=3$
\end{frame}
\fi

% 题库编号：LOG-040
\begin{campuspractice}{题库 LOG-040}
将下式互化：\par
$\log_216=4$
\end{campuspractice}

\ifdefined\TeacherEdition
\begin{frame}[t]{LOG-040：参考解答}
$2^4=16$
\end{frame}
\fi

% 题库编号：LOG-041
\begin{campuspractice}{题库 LOG-041}
将下式互化：\par
$10^{-2}=0.01$
\end{campuspractice}

\ifdefined\TeacherEdition
\begin{frame}[t]{LOG-041：参考解答}
$\lg0.01=-2$
\end{frame}
\fi

% 题库编号：LOG-042
\begin{campuspractice}{题库 LOG-042}
将下式互化：\par
$\log_{\sqrt5}125=6$
\end{campuspractice}

\ifdefined\TeacherEdition
\begin{frame}[t]{LOG-042：参考解答}
$(\sqrt5)^6=125$
\end{frame}
\fi

% 题库编号：LOG-043
\begin{campuspractice}{题库 LOG-043}
以下哪些是正确的直接互化？

A. $10^0=1$ 与 $\lg1=0$； B. $27^{\frac{-1}{3}}=\frac{1}{3}$ 与 $\log_{27}(\frac{1}{3})=\frac{-1}{3}$；

C. $\log_39=2$ 与 $9^{\frac{1}{2}}=3$； D. $\log_55=1$ 与 $5^1=5$。
\end{campuspractice}

\ifdefined\TeacherEdition
\begin{frame}[t]{LOG-043：参考解答}
A、B、D。C两式虽各自成立，但不是底数不变的直接互化。
\end{frame}
\fi

% 题库编号：LOG-044
\begin{campuspractice}{题库 LOG-044}
求下式中的 $x$：\par
$\lg0.01=x$
\end{campuspractice}

\ifdefined\TeacherEdition
\begin{frame}[t]{LOG-044：参考解答}
$x=-2$
\end{frame}
\fi

% 题库编号：LOG-045
\begin{campuspractice}{题库 LOG-045}
求下式中的 $x$：\par
$\log_7(x+2)=2$
\end{campuspractice}

\ifdefined\TeacherEdition
\begin{frame}[t]{LOG-045：参考解答}
$x=47$
\end{frame}
\fi

% 题库编号：LOG-046
\begin{campuspractice}{题库 LOG-046}
求下式中的 $x$：\par
$\log_{\frac{2}{3}}(\frac{9}{4})=x$
\end{campuspractice}

\ifdefined\TeacherEdition
\begin{frame}[t]{LOG-046：参考解答}
$x=-2$
\end{frame}
\fi

% 题库编号：LOG-047
\begin{campuspractice}{题库 LOG-047}
求下式中的 $x$：\par
$x=\log_{\frac{1}{2}}32$
\end{campuspractice}

\ifdefined\TeacherEdition
\begin{frame}[t]{LOG-047：参考解答}
$x=-5$
\end{frame}
\fi

% 题库编号：LOG-048
\begin{campuspractice}{题库 LOG-048}
求下式中的 $x$：\par
$\log_5(x^2)=2$
\end{campuspractice}

\ifdefined\TeacherEdition
\begin{frame}[t]{LOG-048：参考解答}
$x=\pm5$
\end{frame}
\fi

% 题库编号：LOG-049
\begin{campuspractice}{题库 LOG-049}
求下式中的 $x$：\par
$2\log_3x=4$
\end{campuspractice}

\ifdefined\TeacherEdition
\begin{frame}[t]{LOG-049：参考解答}
$x=9$
\end{frame}
\fi

% 题库编号：LOG-050
\begin{campuspractice}{题库 LOG-050}
计算 $2^{-\log_23}$。
\end{campuspractice}

\ifdefined\TeacherEdition
\begin{frame}[t]{LOG-050：参考解答}
$1/3$
\end{frame}
\fi

% 题库编号：LOG-051
\begin{campuspractice}{题库 LOG-051}
计算 $e^{3\ln7}$。
\end{campuspractice}

\ifdefined\TeacherEdition
\begin{frame}[t]{LOG-051：参考解答}
$343$
\end{frame}
\fi

% 题库编号：LOG-052
\begin{campuspractice}{题库 LOG-052}
计算 $\lg(0.001^2)$。
\end{campuspractice}

\ifdefined\TeacherEdition
\begin{frame}[t]{LOG-052：参考解答}
$-6$
\end{frame}
\fi

% 题库编号：LOG-053
\begin{campuspractice}{题库 LOG-053}
求 $x$：$\log_3(\lg x)=1$。
\end{campuspractice}

\ifdefined\TeacherEdition
\begin{frame}[t]{LOG-053：参考解答}
$x=1000$
\end{frame}
\fi

% 题库编号：LOG-054
\begin{campuspractice}{题库 LOG-054}
求 $x$：$\log_3[\log_4(\log_5x)]=0$。
\end{campuspractice}

\ifdefined\TeacherEdition
\begin{frame}[t]{LOG-054：参考解答}
$x=625$
\end{frame}
\fi

% 题库编号：LOG-055
\begin{campuspractice}{题库 LOG-055}
记 $L(x)=\log_3[\log_4(\log_5x)]$。若 $L(x)=1$，求 $x$。
\end{campuspractice}

\ifdefined\TeacherEdition
\begin{frame}[t]{LOG-055：参考解答}
$x=5^{64}$
\end{frame}
\fi

% 题库编号：LOG-056
\begin{campuspractice}{题库 LOG-056}
记 $L(x)=\log_3[\log_4(\log_5x)]$。若 $L(x)=0$，求 $625^{\log_x3}$。
\end{campuspractice}

\ifdefined\TeacherEdition
\begin{frame}[t]{LOG-056：参考解答}
$3$
\end{frame}
\fi

% 题库编号：LOG-057
\begin{campuspractice}{题库 LOG-057}
若 $3^{\log_3[\log_4(\log_5x)]}=1$，求 $x$。
\end{campuspractice}

\ifdefined\TeacherEdition
\begin{frame}[t]{LOG-057：参考解答}
$x=625$
\end{frame}
\fi

% 题库编号：LOG-058
\begin{campuspractice}{题库 LOG-058}
$\log_x16=2$，则 $x$ 等于（\quad）。\par A. $\pm4$\quad B. $4$\quad C. $256$\quad D. $2$
\end{campuspractice}

\ifdefined\TeacherEdition
\begin{frame}[t]{LOG-058：参考解答}
B，$x=4$
\end{frame}
\fi

% 题库编号：LOG-059
\begin{campuspractice}{题库 LOG-059}
若 $f(2^x)=x$（$x\in\mathbb R$），则 $f(5)$ 为（\quad）。\par A. $2^5$\quad B. $5^2$\quad C. $\log_52$\quad D. $\log_25$
\end{campuspractice}

\ifdefined\TeacherEdition
\begin{frame}[t]{LOG-059：参考解答}
D，$f(5)=\log_25$
\end{frame}
\fi

% 题库编号：LOG-060
\begin{campuspractice}{题库 LOG-060}
$\log_3(a+1)=1$，求 $\log_a2+\log_2(a-1)$。
\end{campuspractice}

\ifdefined\TeacherEdition
\begin{frame}[t]{LOG-060：参考解答}
$1$
\end{frame}
\fi

% 题库编号：LOG-061
\begin{campuspractice}{题库 LOG-061}
$a=\log_34$，求 $3^{2a}$。
\end{campuspractice}

\ifdefined\TeacherEdition
\begin{frame}[t]{LOG-061：参考解答}
$16$
\end{frame}
\fi

% 题库编号：LOG-062
\begin{campuspractice}{题库 LOG-062}
将下式在指数式与对数式之间互化：\par
$\log_216=4$
\end{campuspractice}

\ifdefined\TeacherEdition
\begin{frame}[t]{LOG-062：参考解答}
$2^4=16$
\end{frame}
\fi

% 题库编号：LOG-063
\begin{campuspractice}{题库 LOG-063}
将下式在指数式与对数式之间互化：\par
$\log_{\frac{1}{3}}27=-3$
\end{campuspractice}

\ifdefined\TeacherEdition
\begin{frame}[t]{LOG-063：参考解答}
$(1/3)^{-3}=27$
\end{frame}
\fi

% 题库编号：LOG-064
\begin{campuspractice}{题库 LOG-064}
将下式在指数式与对数式之间互化：\par
$\log_{\sqrt3}x=3$
\end{campuspractice}

\ifdefined\TeacherEdition
\begin{frame}[t]{LOG-064：参考解答}
$(\sqrt3)^3=x$
\end{frame}
\fi

% 题库编号：LOG-065
\begin{campuspractice}{题库 LOG-065}
将下式在指数式与对数式之间互化：\par
$5^3=125$
\end{campuspractice}

\ifdefined\TeacherEdition
\begin{frame}[t]{LOG-065：参考解答}
$\log_5125=3$
\end{frame}
\fi

% 题库编号：LOG-066
\begin{campuspractice}{题库 LOG-066}
将下式在指数式与对数式之间互化：\par
$2^{-1}=\frac{1}{2}$
\end{campuspractice}

\ifdefined\TeacherEdition
\begin{frame}[t]{LOG-066：参考解答}
$\log_2(1/2)=-1$
\end{frame}
\fi

% 题库编号：LOG-067
\begin{campuspractice}{题库 LOG-067}
将下式在指数式与对数式之间互化：\par
$(\frac{1}{3})^{-2}=9$
\end{campuspractice}

\ifdefined\TeacherEdition
\begin{frame}[t]{LOG-067：参考解答}
$\log_{1/3}9=-2$
\end{frame}
\fi

% 题库编号：LOG-068
\begin{campuspractice}{题库 LOG-068}
求下式中的 $x$：\par
$\log_x27=\frac{3}{2}$
\end{campuspractice}

\ifdefined\TeacherEdition
\begin{frame}[t]{LOG-068：参考解答}
$x=9$
\end{frame}
\fi

% 题库编号：LOG-069
\begin{campuspractice}{题库 LOG-069}
求下式中的 $x$：\par
$\log_2x=\frac{-2}{3}$
\end{campuspractice}

\ifdefined\TeacherEdition
\begin{frame}[t]{LOG-069：参考解答}
$x=2^{-2/3}$
\end{frame}
\fi

% 题库编号：LOG-070
\begin{campuspractice}{题库 LOG-070}
求下式中的 $x$：\par
$\log_x(3+2\sqrt2)=-2$
\end{campuspractice}

\ifdefined\TeacherEdition
\begin{frame}[t]{LOG-070：参考解答}
$x=\sqrt2-1$
\end{frame}
\fi

% 题库编号：LOG-071
\begin{campuspractice}{题库 LOG-071}
求下式中的 $x$：\par
$\log_5(\log_2x)=0$
\end{campuspractice}

\ifdefined\TeacherEdition
\begin{frame}[t]{LOG-071：参考解答}
$x=2$
\end{frame}
\fi

% 题库编号：LOG-072
\begin{campuspractice}{题库 LOG-072}
求下式中的 $x$：\par
$\log_2[\log_3(\log_2x)]=1$
\end{campuspractice}

\ifdefined\TeacherEdition
\begin{frame}[t]{LOG-072：参考解答}
$x=512$
\end{frame}
\fi

% 题库编号：LOG-073
\begin{campuspractice}{题库 LOG-073}
$a=\log_310,b=\log_37$，求 $3^{a-b}$。\par A. $\frac{10}{7}$\quad B. $\frac{7}{10}$\quad C. $\frac{10}{49}$\quad D. $\frac{49}{10}$
\end{campuspractice}

\ifdefined\TeacherEdition
\begin{frame}[t]{LOG-073：参考解答}
A，$10/7$
\end{frame}
\fi

% 题库编号：LOG-074
\begin{campuspractice}{题库 LOG-074}
$\log_{\sqrt2+1}(3-2\sqrt2)=t$，求 $t$。\par A. $-2$\quad B. $-4$\quad C. $2$\quad D. $4$
\end{campuspractice}

\ifdefined\TeacherEdition
\begin{frame}[t]{LOG-074：参考解答}
A，$-2$
\end{frame}
\fi

% 题库编号：LOG-075
\begin{campuspractice}{题库 LOG-075}
$\log_3(\log_5a)=\log_4(\log_5b)=0$，求 $a/b$。\par A. $1$\quad B. $-1$\quad C. $5$\quad D. $\frac{1}{5}$
\end{campuspractice}

\ifdefined\TeacherEdition
\begin{frame}[t]{LOG-075：参考解答}
A，$a=b=5$，$a/b=1$
\end{frame}
\fi

% 题库编号：LOG-076
\begin{campuspractice}{题库 LOG-076}
解方程 $3^{\log_2x}=\frac{1}{27}$。
\end{campuspractice}

\ifdefined\TeacherEdition
\begin{frame}[t]{LOG-076：参考解答}
$x=1/8$
\end{frame}
\fi

% 题库编号：LOG-077
\begin{campuspractice}{题库 LOG-077}
若 $\log_{1-x}((1+x)^2)=1$，求 $x$。
\end{campuspractice}

\ifdefined\TeacherEdition
\begin{frame}[t]{LOG-077：参考解答}
$x=-3$（剔除底数为1的$x=0$）
\end{frame}
\fi

% 题库编号：LOG-078
\begin{campuspractice}{题库 LOG-078}
计算 $2^{2+\log_23}+3^{2-\log_39}$。
\end{campuspractice}

\ifdefined\TeacherEdition
\begin{frame}[t]{LOG-078：参考解答}
$13$
\end{frame}
\fi

% 题库编号：LOG-079
\begin{campuspractice}{题库 LOG-079}
若 $x=\log_43$，求 $(2^x-2^{-x})^2$。
\end{campuspractice}

\ifdefined\TeacherEdition
\begin{frame}[t]{LOG-079：参考解答}
$4/3$
\end{frame}
\fi

\begin{frame}[t]{数学文化与拓展材料}
以下材料供备课选用，不计入40分钟主课。
\end{frame}

% 来源：图片838—840教材数学文化；按材料概述，不将现代log定义等同早期对数
\begin{frame}[t]{附录：对数为何被发明}
对数曾是简化繁重计算的重要工具。所附教材介绍：
\begin{itemize}
\item 纳皮尔于1614年发表对数著作。
\item 别尔基也独立研究了对数，相关著作于1620年出版。
\item 布里格斯推动常用对数表的编制。
\end{itemize}
\[2^6\cdot2^8=2^{6+8}=2^{14},\qquad64\times256=16384.\]
将幂对应到指数，就能把乘法转化为加法。\end{frame}

% 来源：图片838—839
\begin{frame}[t]{附录：从稀疏表到精细表}
\small
\centering\begin{tabular}{c|rrrrrrrr}
$n$&0&1&2&3&4&5&6&7\\\hline
$2^n$&1&2&4&8&16&32&64&128\\\hline
$n$&8&9&10&11&12&13&14&\\\hline
$2^n$&256&512&1024&2048&4096&8192&16384&
\end{tabular}\par\normalsize
表中缺少3、5、7等数，怎样把间隔缩小？

取接近1的 $q=1.0001$：每次把 $x$ 增加其万分之一，
\[x_0=1,\quad x_{n+1}=x_n+\frac{x_n}{10000}=1.0001x_n,\quad x_n=1.0001^n.\]\end{frame}

% 来源：图片839
\begin{frame}[t]{附录：历史表的近似数}
\[x=1.0001^y,\qquad y=\log_{1.0001}x.\]
\[1.0001^{100}\approx1.010049662,\quad
1.0001^{101}\approx1.010150667.\]
\[1.0001^{46054}\approx99.999955936.\]

教材用接近1的底数解释早期制表思想。数值表经舍入，只能得到近似计算结果；这些现代符号是对思想的解释。\end{frame}

% 来源：图片839—840
\begin{frame}[t]{附录：查表把除法变成减法}
\small\centering\renewcommand{\arraystretch}{1.35}
\begin{tabular}{c|rrrrr}
$y$&20333&20334&20335&20336&20337\\
$x\approx$&7.6385&7.6392&7.6400&7.6408&7.6415\\\hline
$y$&10237&10238&10239&10240&10241\\
$x\approx$&2.7833&2.7836&2.7839&2.7842&2.7844
\end{tabular}\par\normalsize
求 $7640\div2784=7.640\div2.784$。

按表取对应指数 $20335$ 和 $10239$，相减为 $10096$。\end{frame}

% 来源：图片840；修正原近似分母使用精确等号的问题
\begin{frame}[t]{附录：近似计算与换底}
\small\centering\begin{tabular}{c|rrrrr}
$y$&10094&10095&10096&10097&10098\\\hline
$x\approx$&2.7438&2.7441&2.7444&2.7446&2.7449
\end{tabular}\par\normalsize
查表得到商约为 $2.7444$；实际商为 $2.74425287\ldots$。
两者保留4位有效数字均为 $2.744$。

\[\lg x=\frac{\log_{1.0001}x}{\log_{1.0001}10}
\approx\frac{\log_{1.0001}x}{23027}.\]
本页换底公式作为后续学习的预告，分母 $23027$ 是近似值。\end{frame}

% 来源：图片841顶部
\begin{frame}[t]{补充：教材思考题}
由 $0.568^x=0.01$，可准确记为
\[x=\log_{0.568}0.01.\]
借助计算器得 $x\approx8.142$。

对数符号能够准确表示指数，近似小数方便实际计算，两者作用不同。\end{frame}
% 来源：图片841顶部；图片839的制表背景。
\begin{frame}[t]{补充：指数结构与制表背景}
\[2^{\log_23-2}=\frac{2^{\log_23}}{2^2}=\frac34.\]
\[2^{\log_2(3+2)}=5.\]
括号内的加法与指数上的减法，处理方式不同。

所附数学文化材料还介绍：纳皮尔采用接近1的小于1的数
$b=0.9999999$，别尔基采用 $a=1.0001$，以缩小表中相邻数的间隔。

现代指数记号和对数与指数的统一认识经历了长期发展，早期对数表不能直接等同于今天的常用对数表。
\end{frame}

\campusclosing[谢谢]\end{document}
